\subsection{交换群的情形}

在本小节中，我们假设群 \(\mathrm{G}\) 是交换的。

由舒尔引理（VIII, 第 48 页，推论 1），G 的每个单表示的维数都是 1。设 \((\mathrm{M},\pi)\) 是这样一个表示，\(\chi\) 是它的特征标；对每个 \(g\in G\) 与 \(x\in M\)，有 \(\pi(g)(x)=\chi(g)x\)。于是，特征标 \(\chi\) 是从 G 到 K 的乘法群 \(K^{*}\) 的一个同态。反之，从 G 到 \(K^{*}\) 的每个同态都是 G 的一个次数为 1 的表示的特征标。因此，集合 \(\widehat{G}\)（单 K{[}G{]}-模的类所成的集合）可与集合 \(\operatorname{Hom}(G,K^{*})\)（从 G 到 \(K^{*}\) 的同态所成的集合）等同。由此我们在 \(\widehat{G}\) 上导出一个交换群结构；\(\widehat{G}\) 中的乘积对应于表示的张量积。由 VIII, 第 411 页的命题 5，群 G 与 \(\widehat{G}\) 有相同的基数。G 上的每个函数都是中心函数，且 \(\widehat{G}\) 是从 G 到 K 的映射所成的向量空间的一组基（同前引文）。由特征标的正交关系，这样的映射 f 关于基 \(\widehat{G}\) 有如下唯一的分解：

\[
f = \sum_ {\chi \in \widehat {\mathrm{G}}} \langle \chi , f \rangle_ {\mathrm{G}} \chi .\tag{39}
\]

对映射 \(f\) 与 \(f'\)（从 \(G\) 到 \(K\) 的映射），我们有关系式

\[
\langle f, f ^ {\prime} \rangle_ {\mathrm{G}} = \sum_ {\chi \in \widehat {\mathrm{G}}} \langle \chi , f \rangle_ {\mathrm{G}} \langle \chi , f ^ {\prime} \rangle_ {\mathrm{G}}.\tag{40}
\]

设 \((\mathrm{V},\pi)\) 是 G 的一个线性表示。对任意 \(\chi\in\widehat{\mathrm{G}}\)，用 \(V^{\chi}\) 表示 V 中由满足 \(\pi(g)(v)=\chi(g)v\)（对每个 \(g\in\mathrm{G}\)）的向量 v 组成的子空间。空间 \(V^{\chi}\) 是 \(\chi\) 型同型分量（K{[}G{]}-模 V 的同型分量）。空间 V 是族 \((V^{\chi})_{\chi\in\widehat{\mathrm{G}}}\) 的直和，而与这一分解相伴的 V 的投影算子 \(p_{\chi}\)（以 \(V^{\chi}\) 为像的投影算子）由

\[
p _ {\chi} = | \mathrm{G} | ^ {- 1} \sum_ {g \in \mathrm{G}} \chi (g ^ {- 1}) \pi (g)\tag{41}
\]

给出（依关系式 (19)）。

注. —— 设 n 是群 G 的基数，\(\mu_{n}(\mathrm{K})\) 是 K 中 n 次单位根所成的群。对每个 \(g \in G\)，有 \(g^{n} = 1\)；于是，\(\widehat{G}\) 可与群 \(\operatorname{Hom}(\mathbf{G}, \mu_{n}(\mathbf{K}))\) 等同。群 \(\mu_{n}(\mathbf{K})\) 是 n 阶循环群（V, §11, No.~2, 第 78 页，定理 1）。因此群 \(\widehat{G}\) 同构于群 \(\mathrm{D}(\mathrm{G}) = \operatorname{Hom}(\mathrm{G}, \mathbf{Q}/\mathbf{Z})\)。由 VII, §4, No.~9, 第 26 页的

命题 10，群 \(\widehat{\mathbf{G}}\) 同构于群 \(\mathbf{G}\)，且把元素 \(g\)（\(\mathbf{G}\) 的元素）映为同态 \(\chi \mapsto \chi(g)\)（从 \(\widehat{\mathbf{G}}\) 到 \(\mathrm{K}^*\) 的同态）的映射是从 \(\mathbf{G}\) 到 \(\widehat{\widehat{\mathbf{G}}}\) 的一个同构。
% semantic-fragments: legacy-input-seam
\relax{}
